(1) 令 $\Omega_R=\sqrt{\Omega^2+\Delta^2}$。因无迹矩阵 $H^2=(\hbar^2\Omega_R^2/4)\mathsf I$，本征能量为 $E_\pm=\pm\hbar\Omega_R/2$，能量劈裂 $\hbar\Omega_R$。这是在驱动旋转框架内的准能量结构，裸能级间隔、驱动频率与失谐的物理定义要一致，不能把此间隔误认为无驱动时的原子光谱间隔。
(2) 由于 $H^2$ 与单位矩阵成正比，演化算符可精确写为 $U(t)=\cos(\Omega_Rt/2)\mathsf I-2i\sin(\Omega_Rt/2)H/(\hbar\Omega_R)$。初始态 $|g\rangle$ 演化到 $|e\rangle$ 的振幅为 $-i(\Omega/\Omega_R)\sin(\Omega_Rt/2)$，所以 $P_e(t)=\Omega^2/(\Omega^2+\Delta^2)\sin^2[\sqrt{\Omega^2+\Delta^2}\,t/2]$。
(3) 在精确共振 $\Delta=0$ 时 $\Omega_R=\Omega$，有 $P_e(t)=\sin^2(\Omega t/2)$，首次达到一在 $t_\pi=\pi/\Omega$，称为 $\pi$ 脉冲。对于固定非零失谐 $\Delta$，即使选择最佳脉冲长度，最大概率也只有 $P_{e,\max}=\Omega^2/(\Omega^2+\Delta^2)<1$，振荡周期变短但幅度降低，不能通过简单延长单个矩形脉冲弥补失谐。
(4) 对连续相干驱动、观察时间包含大量 Rabi 周期，正弦平方的时间平均为一半，故 $\overline P_e=\Omega^2/[2(\Omega^2+\Delta^2)]$。这些结论以驱动近单色、$\Omega$ 与 $|\Delta|$ 远小于光学载频及二能级近似成立为前提；实际退相干与自发辐射使长期振荡振幅衰减，平均占据需由含弛豫的主方程重新计算。
